% Reference draft, mathematics companion -- MH370 Bayesian reconstruction
% INTERNAL. Not the paper. Generated by build_reference_draft.py -- do not hand-edit.
% Every equation is copied from the project note named above it. Where a note gives no
% equation, none appears here.
\documentclass[11pt,a4paper]{article}
\usepackage[utf8]{inputenc}\usepackage{amsmath,amssymb,geometry,fancyhdr}
\geometry{margin=2.54cm}
\pagestyle{fancy}\fancyhf{}\fancyhead[L]{Reference draft -- internal, not the paper}
\fancyhead[R]{\thepage}
\title{MH370 Bayesian reconstruction:\\ governing equations of the reference draft}
\author{Internal reference --- paper drafting and citations module}
\date{12 October 2026}
\begin{document}\maketitle
\noindent\textbf{Internal document.} This is the mathematics companion to
\texttt{paper/REFERENCE-DRAFT.md}. Equations are reproduced as the project notes give
them, each with the file that implements it. Nothing here is a derivation by the
drafting module, and nothing is included that a note does not state.
\medskip
\section{Reproduction}
\subsection{OU stationary initial deviations (Mach, angle, wind)}
\noindent\textit{Status:} BUILT. \textit{Source note:} \texttt{results/davey-2016-reference.md}.
\medskip
\begin{equation}
\sigma = \sqrt{q / 2\beta}
\end{equation}
\noindent\textbf{Symbols.}
\begin{verbatim}
\sigma = stationary standard deviation of the OU process [Mach, degrees, or knots
depending on state]; q = OU process noise intensity [state-dependent units]; \beta =
OU reversion rate [s^{-1}]
\end{verbatim}
\noindent\textbf{Implemented in.}
\begin{verbatim}
crates/flight/src/lib.rs
\end{verbatim}
\medskip
\subsection{Mach acceleration rate}
\noindent\textit{Status:} BUILT. \textit{Source note:} \texttt{results/davey-2016-reference.md}.
\medskip
\begin{equation}
mach\_rate\_per\_s = 0.1/60
\end{equation}
\noindent\textbf{Symbols.}
\begin{verbatim}
mach\_rate\_per\_s = assumed maximum rate of Mach number change [Mach/s]
\end{verbatim}
\medskip
\subsection{Davey ch.11 §11.2 single-flaperon drift likelihood update (GDP drifter reproduction)}
\noindent\textit{Status:} MIXED. \textit{Source note:} \texttt{results/davey-ch11-reproduction.md}.
\medskip
\begin{equation}
p(x_{final} \mid Z_K) \approx \sum_p w_p \delta(x_{final} - x_p)
\end{equation}
\noindent\textbf{Symbols.}
\begin{verbatim}
x_{final} = final (impact) position; Z_K = observation set up to time K; w_p = weight
of particle p; x_p = position of particle p; \delta = Dirac delta
\end{verbatim}
\medskip
\begin{equation}
p(x_{final} \mid Z_K, y) \propto p(y \mid x_{final}) p(x_{final} \mid Z_K) \approx \sum_p w_p p(y \mid x_p) \delta(x_{final} - x_p)
\end{equation}
\noindent\textbf{Symbols.}
\begin{verbatim}
y = new observation; p(y \mid x_{final}) = likelihood of y given the final position;
w_p = particle weight; x_p = particle p position
\end{verbatim}
\noindent\textbf{Implemented in.}
\begin{verbatim}
engine/hypotheses/debris-drift/prepare/gdp/davey_ch11.py
\end{verbatim}
\medskip
\begin{equation}
p_{drifter}(x, y) \approx \frac{1}{P} \sum_p \delta(x - x_p) \delta(y - y_p)
\end{equation}
\noindent\textbf{Symbols.}
\begin{verbatim}
x = drifter location on 8 March 2014; y = drifter location on 29 July 2015 (508 days
later); P = number of drifter trajectories; x_p, y_p = p-th trajectory's start/end
positions
\end{verbatim}
\noindent\textbf{Implemented in.}
\begin{verbatim}
engine/hypotheses/debris-drift/prepare/gdp/davey_ch11.py
\end{verbatim}
\medskip
\begin{equation}
l(x_{final}) = \int_R p(y \mid x_{final})\, dy
\end{equation}
\noindent\textbf{Symbols.}
\begin{verbatim}
R = small region surrounding Reunion Island, size not printed; x_{final} = final
position; y = arrival location
\end{verbatim}
\noindent\textbf{Implemented in.}
\begin{verbatim}
engine/hypotheses/debris-drift/prepare/gdp/davey_ch11.py
\end{verbatim}
\medskip
\begin{equation}
l(x) \approx \frac{\frac{1}{P}\sum_p K(x - x_p) + \epsilon}{\frac{1}{\tilde P}\sum_p K(x - \tilde x_p) + \epsilon}
\end{equation}
\noindent\textbf{Symbols.}
\begin{verbatim}
K = fixed kernel, 1 degree standard deviation in latitude and longitude; \epsilon =
floor constant, 10^{-4} as printed; x_p = trajectories passing through R in the right
time window, joined ones included; \tilde x_p = overall set of drifter locations,
February-April; P, \tilde P = respective sample counts
\end{verbatim}
\noindent\textbf{Implemented in.}
\begin{verbatim}
engine/hypotheses/debris-drift/prepare/gdp/davey_ch11.py
\end{verbatim}
\medskip
\begin{equation}
\text{ESS fraction} = \frac{(\sum w l)^2}{(\sum w \cdot \sum w l^2)}
\end{equation}
\noindent\textbf{Symbols.}
\begin{verbatim}
w = particle weight; l = relative likelihood at that particle
\end{verbatim}
\medskip
\begin{equation}
w = \text{density} \times 0.0625
\end{equation}
\noindent\textbf{Symbols.}
\begin{verbatim}
w = particle weight; density = per-deg^2 density at the cell centre; 0.0625 = area in
deg^2 of a 0.25 degree cell
\end{verbatim}
\noindent\textbf{Implemented in.}
\begin{verbatim}
summary.rs::map_cell
\end{verbatim}
\medskip
\begin{equation}
l / E_{prior}[l]
\end{equation}
\noindent\textbf{Symbols.}
\begin{verbatim}
l = relative likelihood; E_{prior}[l] = expectation of l under the prior
\end{verbatim}
\medskip
\section{Update}
\subsection{Composer pass 1 (run C) — Bayesian evidence/position composition pipeline}
\noindent\textit{Status:} MIXED. \textit{Source note:} \texttt{results/composer-pass1-run-c.md}.
\medskip
\begin{equation}
\ln R = \ln E[L_P \mid \text{flight} + \text{EoF} + \text{drift} + \text{hydro}] - \ln E[L_P \mid \text{flight} + \text{EoF}]
\end{equation}
\noindent\textbf{Symbols.}
\begin{verbatim}
R = evidence ratio of H (Pléiades-conditioned) against G [dimensionless]; L_P =
per-row Pléiades likelihood [dimensionless]; E[\cdot] = expectation over composed
impact rows [dimensionless]
\end{verbatim}
\medskip
\begin{equation}
Ga = G \times L_{search}
\end{equation}
\noindent\textbf{Symbols.}
\begin{verbatim}
Ga = product G after applying searched-area non-detection; G = composed
flight+EoF+drift+hydro product; L_{search} = per-row point-target non-detection
likelihood [dimensionless]
\end{verbatim}
\medskip
\begin{equation}
w = W_Q(\text{stratum}) \times w_{pooled}
\end{equation}
\noindent\textbf{Symbols.}
\begin{verbatim}
w = composed per-row weight written to the exchange [dimensionless];
W_Q(\text{stratum}) = stratum mixture weight [dimensionless]; w_{pooled} = pooled
within-stratum weight across the 8 replicate seeds [dimensionless]
\end{verbatim}
\noindent\textbf{Implemented in.}
\begin{verbatim}
mh370-exchange/composer/next-run-c/row-weights/<option>-<product>/
\end{verbatim}
\medskip
\begin{equation}
H_{r0} - H_{c} = 2\ln A_{scene}
\end{equation}
\noindent\textbf{Symbols.}
\begin{verbatim}
H_{r0} = evidence of product H with COSMO at correlated-error ρ=0; H_{c} = evidence of
product H plus COSMO; A_{scene} = scene-area normalisation constant omitted by the ρ
surfaces [dimensionless]
\end{verbatim}
\medskip
\section{Modification}
\subsection{Resampling (branching scheme)}
\noindent\textit{Status:} MIXED. \textit{Source note:} \texttt{results/davey-2016-reference.md}.
\medskip
\begin{equation}
w' = w / \bar{n}
\end{equation}
\noindent\textbf{Symbols.}
\begin{verbatim}
w' = weight of each branch after duplication [dimensionless]; w = weight of the parent
particle before branching [dimensionless]; \bar{n} = number of branches spawned
[dimensionless]
\end{verbatim}
\medskip
\section{Extension}
\subsection{Fuel and endurance model}
\noindent\textit{Status:} MIXED. \textit{Source note:} \texttt{results/core-model-stages.md}.
\medskip
\begin{equation}
\text{flow} = aM^2 + \frac{b}{M^2}
\end{equation}
\noindent\textbf{Symbols.}
\begin{verbatim}
a = fitted drag-law coefficient [flow units]; b = fitted drag-law coefficient [flow
units]; M = Mach number [dimensionless]; flow = fuel flow rate [kg/h per engine]
\end{verbatim}
\noindent\textbf{Implemented in.}
\begin{verbatim}
crates/flight/src/fuel.rs
\end{verbatim}
\medskip
\begin{equation}
36{,}609 = 43{,}800 - 0.700 \times (43{,}800 - 33{,}524)
\end{equation}
\noindent\textbf{Symbols.}
\begin{verbatim}
43{,}800 = ACARS fuel reading at 17:06:43 [kg]; 33{,}524 = Boeing's computed fuel at
arc 1 [kg]; 0.700 = fraction of the 78.73 min segment elapsed by 18:01:49
[dimensionless]; 36{,}609 = initial fuel at filter start 18:01:49 [kg]
\end{verbatim}
\noindent\textbf{Implemented in.}
\begin{verbatim}
crates/flight/src/fuel.rs
\end{verbatim}
\medskip
\subsection{Drifting BFO bias}
\noindent\textit{Status:} MIXED. \textit{Source note:} \texttt{results/core-model-stages.md}.
\medskip
\begin{equation}
\text{rate} = \frac{7^2 - 0.995^2}{3600} = 0.01334 \text{ Hz}^2/\text{s}
\end{equation}
\noindent\textbf{Symbols.}
\begin{verbatim}
7 = book's measurement standard deviation [Hz]; 0.995 = genuinely random part of the
measurement sd [Hz]; 3600 = hourly handshake spacing [s]; rate = bias random-walk
variance growth rate [Hz^2/s]
\end{verbatim}
\medskip
\begin{equation}
\text{variance} \mathrel{+}= \text{rate} \times \text{gap}
\end{equation}
\noindent\textbf{Symbols.}
\begin{verbatim}
variance = bias random-walk variance [Hz^2]; rate = drift rate [Hz^2/s]; gap = time
between successive BFO epochs [s]
\end{verbatim}
\medskip
\subsection{Debris-drift likelihood term (ocean drift module)}
\noindent\textit{Status:} MIXED. \textit{Source note:} \texttt{results/debris-drift-methods-draft.md}.
\medskip
\begin{equation}
\ln p(D_{\text{drift}} \mid x)
\end{equation}
\noindent\textbf{Symbols.}
\begin{verbatim}
x = impact location; D_drift = recovered-debris record (where and when each identified
item was found)
\end{verbatim}
\noindent\textbf{Implemented in.}
\begin{verbatim}
interpolate.rs
\end{verbatim}
\medskip
\begin{equation}
k_{ij} = \phi(d(y_j, z_i)/h)/h
\end{equation}
\noindent\textbf{Symbols.}
\begin{verbatim}
k_ij = kernel weight linking find j to particle i; d(y_j,z_i) = distance between find
location y_j and particle beaching location z_i [km]; h = kernel bandwidth, 50 km
[km]; \phi = kernel density function
\end{verbatim}
\medskip
\begin{equation}
a_{jb} = \frac{1}{N} \sum_i k_{ij} \cdot P(\text{discovery in } [t_j, t_j+w_j) \cap \text{period } b \mid \tau_i) / w_j
\end{equation}
\noindent\textbf{Symbols.}
\begin{verbatim}
a_jb = block coefficient for find j, block b; N = particles in ensemble; k_ij = kernel
weight; t_j = find time; w_j = width of discovery window [days]; \tau_i = beaching
time of particle i
\end{verbatim}
\medskip
\begin{equation}
A_{cb} = \frac{1}{N} \sum_i m_s(z_i) \cdot P(\text{discovery in period } b \text{ before the window end} \mid \tau_i)
\end{equation}
\noindent\textbf{Symbols.}
\begin{verbatim}
A_cb = class c's recovery mass in block b; m_s(z_i) = membership of particle i's
beaching location in segment s; \tau_i = beaching time
\end{verbatim}
\medskip
\begin{equation}
\ln \nu_b \sim N(\mu_{\text{period}}, 1)
\end{equation}
\noindent\textbf{Symbols.}
\begin{verbatim}
\nu_b = latent identification level for block b; \mu_period = period mean, ln 0.1, ln
0.5, ln 1 for I1-I3
\end{verbatim}
\medskip
\begin{equation}
L(x) = E_\eta E_\nu \prod_j \left[ \frac{\sum_b \nu_b a_{jb}(x)}{\sum_b \nu_b A_{c(j)b}(x)} \right]
\end{equation}
\noindent\textbf{Symbols.}
\begin{verbatim}
L(x) = likelihood at source node x; E_\eta = expectation over environment draws; E_\nu
= expectation over identification-level draws; a_jb(x) = block coefficient; A_c(j)b(x)
= recovery mass for find j's class
\end{verbatim}
\medskip
\begin{equation}
\mathrm{var}(\ln L_A - \ln L_B)/4
\end{equation}
\noindent\textbf{Symbols.}
\begin{verbatim}
L_A, L_B = likelihood estimates from the two split halves
\end{verbatim}
\medskip
\begin{equation}
\sigma_{\mathrm{eff}}^2 = \sigma^2 - K_{\mathrm{ref}}/T
\end{equation}
\noindent\textbf{Symbols.}
\begin{verbatim}
\sigma_eff = effective error-field standard deviation; \sigma = GDP drifter-replay
error, about 0.09-0.12 m/s [m/s]; K_ref = reference diffusivity, 248 m^2/s; T =
decorrelation time, about 5-15 days [days]
\end{verbatim}
\medskip
\subsection{Hydroacoustics module (SOFAR + AGW branches, detection/likelihood gating)}
\noindent\textit{Status:} MIXED. \textit{Source note:} \texttt{results/hydroacoustics-methods-draft.md}.
\medskip
\begin{equation}
E_s = \eta \cdot E \cdot \rho c / 2\pi
\end{equation}
\noindent\textbf{Symbols.}
\begin{verbatim}
E_s = source energy flux density at 1 m [not stated]; \eta = acoustic coupling
efficiency [dimensionless]; E = impact kinetic energy [J]; \rho = water density
[kg/m^3]; c = sound speed [m/s]
\end{verbatim}
\medskip
\begin{equation}
\eta_{cal} = (0.7^2/R) / [E \rho c/2\pi \cdot \Sigma_b S_b(\tau) 10^{(-TL_b/10)}]
\end{equation}
\noindent\textbf{Symbols.}
\begin{verbatim}
\eta_{cal} = calibrated coupling efficiency [dimensionless]; 0.7 = observed peak
pressure at H11 [Pa]; R = term in denominator (meaning not stated in note); E = impact
kinetic energy [J]; \rho c = acoustic impedance; S_b(\tau) = source spectrum in band
b; TL_b = transmission loss in band b [dB]
\end{verbatim}
\noindent\textbf{Implemented in.}
\begin{verbatim}
hydroacoustics-f35a-eta-calibration.md
\end{verbatim}
\medskip
\begin{equation}
T_c \approx 4H/c
\end{equation}
\noindent\textbf{Symbols.}
\begin{verbatim}
T_c = AGW cut-off period [s]; H = local water depth [m]; c = sound speed [m/s]
\end{verbatim}
\noindent\textbf{Implemented in.}
\begin{verbatim}
agw_regime.py
\end{verbatim}
\medskip
\begin{equation}
J = m v_z
\end{equation}
\noindent\textbf{Symbols.}
\begin{verbatim}
J = vertical surface-force impulse [N\cdot s]; m = impactor mass [kg]; v_z = vertical
impact speed [m/s]
\end{verbatim}
\noindent\textbf{Implemented in.}
\begin{verbatim}
hydroacoustics-agw-scenarios.md
\end{verbatim}
\medskip
\subsection{Pléiades/COSMO-SkyMed conditional hypothesis module (positional likelihood, two-epoch windage calibration, tension statistics)}
\noindent\textit{Status:} MIXED. \textit{Source note:} \texttt{results/pleiades-methods-draft.md}.
\medskip
\begin{equation}
p(y_c \mid s) = \frac{1}{N_w} \sum_k N_2\!\left(y_c ; x_k(s, t_c), \mathrm{diag}(v_E, v_N)\right)
\end{equation}
\noindent\textbf{Symbols.}
\begin{verbatim}
y_c = position of target object or cluster c [deg]; s = candidate impact location; N_w
= number of windage nodes; x_k(s,t_c) = deterministic forward-transport endpoint of
impact s under windage node k at scene time t_c; t_c = Pléiades/COSMO scene
acquisition time [UTC]; v_E = east-component displacement variance [km^2]; v_N =
north-component displacement variance [km^2]
\end{verbatim}
\noindent\textbf{Implemented in.}
\begin{verbatim}
export.rs (pleiades_export_likelihood_surface)
\end{verbatim}
\medskip
\begin{equation}
v_i(\Delta t) = 2\sigma_i^2 T_i\left[\Delta t - T_i\left(1 - e^{-\Delta t / T_i}\right)\right] + sd_{target}^2
\end{equation}
\noindent\textbf{Symbols.}
\begin{verbatim}
v_i = variance of displacement component i (east or north) [km^2]; Δt = elapsed time
from impact to scene acquisition [d]; σ_i = OU-process velocity-error standard
deviation for component i [m/s]; T_i = OU velocity-error correlation time for
component i [d]; sd_target = target position uncertainty [km]
\end{verbatim}
\medskip
\begin{equation}
L(s \mid H) = \sum_c w_c\, p(y_c \mid s) \cdot A_{scene}
\end{equation}
\noindent\textbf{Symbols.}
\begin{verbatim}
w_c = declared weight of cluster/object c; p(y_c|s) = positional likelihood of target
c given impact s; A_scene = scene area [km^2]
\end{verbatim}
\noindent\textbf{Implemented in.}
\begin{verbatim}
pleiades_export_likelihood_surface
\end{verbatim}
\medskip
\begin{equation}
\tilde{d} = 2\,\mathrm{Var}_P\!\left[\log P/\pi\right]
\end{equation}
\noindent\textbf{Symbols.}
\begin{verbatim}
d̃ = Bayesian model dimensionality; P = posterior density; π = prior density
\end{verbatim}
\medskip
\begin{equation}
d = \tilde{d}_A + \tilde{d}_B - \tilde{d}_{AB}
\end{equation}
\noindent\textbf{Symbols.}
\begin{verbatim}
d = shared Bayesian model dimensionality between data sets A and B; d̃_A, d̃_B =
individual dimensionalities; d̃_AB = joint dimensionality
\end{verbatim}
\medskip
\begin{equation}
p = P(\chi^2_d > d - 2\ln S)
\end{equation}
\noindent\textbf{Symbols.}
\begin{verbatim}
p = tension probability between two posteriors; d = shared Bayesian model
dimensionality; S = suspiciousness statistic
\end{verbatim}
\medskip
\begin{equation}
\ln S = \ln R - \ln I
\end{equation}
\noindent\textbf{Symbols.}
\begin{verbatim}
S = suspiciousness; R = Bayes ratio against a flat prior; I = information ratio
\end{verbatim}
\medskip
\begin{equation}
\rho + (1-\rho)\prod_k (1 - q_k c_k)
\end{equation}
\noindent\textbf{Symbols.}
\begin{verbatim}
ρ = shared miss-probability floor (0.05); q_k = search effectiveness of campaign k;
c_k = coverage fraction of campaign k
\end{verbatim}
\noindent\textbf{Implemented in.}
\begin{verbatim}
mh370 evaluate (searched-areas module per-impact column)
\end{verbatim}
\medskip
\subsection{Non-detection likelihood for seabed search campaigns (searched-areas module)}
\noindent\textit{Status:} MIXED. \textit{Source note:} \texttt{results/seabed-search-methods-draft.md}.
\medskip
\begin{equation}
P(no find | y) = \rho + (1 - \rho) \prod_k [1 - c_k(y) g_k(W) q_k]
\end{equation}
\noindent\textbf{Symbols.}
\begin{verbatim}
y = impact position [-]; W = wreckage field produced by the impact [-]; \rho =
probability the field is of the undetectable class [-]; c_k(y) = fraction of ground
around y on which campaign k holds valid data [-]; g_k(W) = size response of campaign
k's sensor to field W [-]; q_k = probability of recognition given detectable data over
the field, campaign k [-]
\end{verbatim}
\medskip
\begin{equation}
p(x\_final | S, Z_K) \propto [1 - P_D(x\_final)] p(x\_final | Z_K)
\end{equation}
\noindent\textbf{Symbols.}
\begin{verbatim}
x_final = final (impact) position [-]; P_D = cumulative probability the search effort
would have detected the aircraft at that location [-]; Z_K = search evidence up to
increment K [-]
\end{verbatim}
\medskip
\begin{equation}
P(no find | y) = 1 - c(y) q = 1 - P_D(y)
\end{equation}
\noindent\textbf{Symbols.}
\begin{verbatim}
c(y) = covered fraction of ground around y for a single campaign [-]; q = probability
of recognition given detectable data [-]; P_D(y) = Davey's detection probability at y
[-]
\end{verbatim}
\noindent\textbf{Implemented in.}
\begin{verbatim}
module test reduces_to_davey_eq_11_1_for_a_point_target_one_campaign_and_rho_zero
\end{verbatim}
\medskip
\begin{equation}
P(find during search of A) = P_D \int_A p(x | Z_K) dx
\end{equation}
\noindent\textbf{Symbols.}
\begin{verbatim}
A = candidate search area [km^2]; P_D = planning (constant) probability of detection
[-]; p(x | Z_K) = residual impact density given evidence to date [-]
\end{verbatim}
\medskip
\subsection{Wreckage settling (ocean-surface-impact-to-seabed transform)}
\noindent\textit{Status:} MIXED. \textit{Source note:} \texttt{results/settling-methods-draft.md}.
\medskip
\begin{equation}
P(\text{intact}) = L(\ln(a/V_d)/s)\, L(\ln(c/V)/s)
\end{equation}
\noindent\textbf{Symbols.}
\begin{verbatim}
L = logistic function [dimensionless]; a = intact-threshold speed, 8 m/s [m/s]; V_d =
descent speed at first contact [m/s]; c = threshold total speed, 100 m/s [m/s]; V =
total speed at first contact [m/s]; s = logistic scale parameter, 0.2 [dimensionless]
\end{verbatim}
\noindent\textbf{Implemented in.}
\begin{verbatim}
breakup.rs
\end{verbatim}
\medskip
\begin{equation}
P(\text{fragmented}) = (1 - P(\text{intact}))\, L(\ln(V/b)/s)
\end{equation}
\noindent\textbf{Symbols.}
\begin{verbatim}
b = fragmentation-threshold speed, 110 m/s [m/s]; V = total speed at first contact
[m/s]; s = logistic scale parameter, 0.2 [dimensionless]
\end{verbatim}
\noindent\textbf{Implemented in.}
\begin{verbatim}
breakup.rs
\end{verbatim}
\medskip
\begin{equation}
P(\text{broken}) = \text{the remainder}
\end{equation}
\noindent\textbf{Symbols.}
\begin{verbatim}
P(\text{broken}) = probability of the broken breakup family [dimensionless]
\end{verbatim}
\noindent\textbf{Implemented in.}
\begin{verbatim}
breakup.rs
\end{verbatim}
\medskip
\begin{equation}
w = k \sqrt{2 g s (1 - \rho_w/\rho_m) / (\rho_w C_d)}
\end{equation}
\noindent\textbf{Symbols.}
\begin{verbatim}
w = descent speed relative to the water [m/s]; k = descent factor for
fluttering/tumbling bodies, k >= 1 [dimensionless]; g = gravitational acceleration
[m/s^2]; s = areal density [kg/m^2]; rho_w = water density [kg/m^3]; rho_m = material
density [kg/m^3]; C_d = drag coefficient [dimensionless]
\end{verbatim}
\medskip
\begin{equation}
v_{ground} = w - w_{up}
\end{equation}
\noindent\textbf{Symbols.}
\begin{verbatim}
v_{ground} = descent speed relative to the ground [m/s]; w = descent speed relative to
water [m/s]; w_{up} = vertical water velocity [m/s]
\end{verbatim}
\medskip
\subsection{Composer end-to-end evidence composition (core run I): combining flight + end-of-flight + drift + searches + Pléiades + hydroacoustics into posterior impact/seabed distributions}
\noindent\textit{Status:} MIXED. \textit{Source note:} \texttt{results/composer-pass2-run-i.md}.
\medskip
\begin{equation}
L = p f_{lag} + (1-p)/W
\end{equation}
\noindent\textbf{Symbols.}
\begin{verbatim}
L = combined cause-model likelihood [dimensionless]; p = mixing weight on cause C1
(fuel exhaustion + APU), p=0.75 under q=1/3 (PRIMARY), q=1 (sensitivity)
[dimensionless]; f_lag = Erlang-lag likelihood component for C1 (unpowered)
[dimensionless]; W = normalising window term for cause C2 (silent) [dimensionless]; q
= parameter setting p (q=1/3 primary, q=1 sensitivity) [dimensionless]
\end{verbatim}
\noindent\textbf{Implemented in.}
\begin{verbatim}
option_table_0019.seed_terms
\end{verbatim}
\medskip
\begin{equation}
L_M = \frac{1}{2}(L_G + L_C),\quad L_m = drift_m \times Pleiades_m
\end{equation}
\noindent\textbf{Symbols.}
\begin{verbatim}
L_M = row likelihood under the equal-weight ocean-model mixture [dimensionless]; L_G =
row likelihood under the GLORYS12 ocean model [dimensionless]; L_C = row likelihood
under the GlobCurrent ocean model [dimensionless]; L_m = per-model likelihood for
ocean model m (drift factor times Pléiades factor) [dimensionless]
\end{verbatim}
\medskip
\begin{equation}
P(GLORYS12 \mid data) = Z_G / (Z_G + Z_C)
\end{equation}
\noindent\textbf{Symbols.}
\begin{verbatim}
P(GLORYS12|data) = posterior probability of the GLORYS12 ocean model [dimensionless];
Z_G = marginal likelihood (evidence) under GLORYS12 [dimensionless]; Z_C = marginal
likelihood under GlobCurrent [dimensionless]
\end{verbatim}
\medskip
\begin{equation}
W_s = P_{core} \times Z_s(00{:}19\ cell)
\end{equation}
\noindent\textbf{Symbols.}
\begin{verbatim}
W_s = stratum mixing weight for core family s under ruling C (PRIMARY)
[dimensionless]; P_core = equal-prior probability of core family s [dimensionless];
Z_s(00:19 cell) = evidence of family s evaluated at the 00:19 data cell
[dimensionless]
\end{verbatim}
\medskip
\begin{equation}
W_s = P_{core} \times Z_s(product)
\end{equation}
\noindent\textbf{Symbols.}
\begin{verbatim}
W_s = stratum mixing weight for family s under the all-evidence mixture
[dimensionless]; P_core = equal-prior probability of family s [dimensionless];
Z_s(product) = evidence of family s evaluated at the full product [dimensionless]
\end{verbatim}
\medskip
\end{document}
